\section{Conclusion}
\label{sec:conclusion}

Executing Retrieval-Augmented Generation on mobile devices is fundamentally constrained by the severe computational cost of prompt prefill on low-power CPUs. While reusing precomputed Key-Value caches can eliminate redundant prefill computation, existing mobile inference runtimes treat KV-cache files as opaque session snapshots, introducing grave vulnerabilities: silent stale-cache hallucinations, unmanaged context-budget overflows, and continuation boundary misalignments.

In this paper, we presented \textbf{HACER-RAG}, an evidence-version-aware inference runtime that makes KV-cache reuse reliable and context-safe on commodity Android smartphones. HACER-RAG pairs a canonical evidence layout with cryptographic chunk-level hashing to bind persistent KV tensors to explicit document versions. Under corpus mutations, HACER-RAG executes a selective repair policy that preserves valid prefix KV tensors while triggering online prefill only for updated downstream tokens. Furthermore, HACER-RAG introduces a JNI-level context-budget safety guard that dynamically enforces headroom constraints and accounts for continuation boundary-token alignment.

We evaluated HACER-RAG on a physical Xiaomi Redmi 12 5G smartphone across 90 matched baseline/reuse pairs and 9 repeated multi-chunk trials. The empirical results demonstrate that HACER-RAG reduces prompt prefill compute latency by up to \textbf{86.67\%}, saving over \textbf{46.9~seconds} of CPU execution and delivering up to \textbf{3.65$\times$} end-to-end wall-clock speedup with 100\% deterministic output parity and zero engine truncations. A balanced Two-Way ANOVA confirmed that compute reductions are strictly governed by cached token count ($F = 1920.56, p < 10^{-68}$), with no statistically significant effect from allocated context capacity ($p = 0.9808$). In repeated multi-chunk trials, HACER-RAG achieved an average prefill compute reduction of \textbf{44.13\%} and a wall speedup of \textbf{1.63$\times$} while preserving 100\% multi-hop factual recall, incurring an average flash cache-load overhead of only 256.0~ms. Across five controlled dynamic corpus scenarios, HACER-RAG eliminated silent stale hallucinations and runtime crashes where naive session reuse failed.

By rigorously characterizing the $>150:1$ operational asymmetry between mobile flash storage bandwidth and CPU prefill throughput, this work establishes that evidence-version-aware KV-cache reuse provides a practical, robust foundation for high-performance, private, and fresh on-device conversational AI.
